% Cover letter to Artificial Intelligence in Medicine (Elsevier). Compile with pdflatex.
\documentclass[11pt,a4paper]{letter}
\usepackage[T1]{fontenc}\usepackage[utf8]{inputenc}\usepackage{lmodern}\usepackage{microtype}
\usepackage[margin=2.5cm]{geometry}\usepackage[hidelinks]{hyperref}
\signature{Dr Aminat Abiola Ajibola\\Assistant Professor\\College of Computer Science and Engineering\\University of Hafr Al-Batin, Eastern Province, Kingdom of Saudi Arabia\\aajibola@uhb.edu.sa\\[0.6em]on behalf of the co-author Roger Nick Anaedevha\\Institute of Cyber Intelligence Systems, National Research Nuclear University MEPhI\\Moscow, Russian Federation\\roger@robustidps.ai}
\address{College of Computer Science and Engineering\\University of Hafr Al-Batin\\Eastern Province, Kingdom of Saudi Arabia\\aajibola@uhb.edu.sa}
\date{\today}
\begin{document}
\begin{letter}{The Editor-in-Chief\\\emph{Artificial Intelligence in Medicine}\\Elsevier}
\opening{Dear Editor,}

We submit the manuscript \textbf{``Machine Learning Models for Mental-Health Prediction from
Actigraphy: A Systematic Appraisal and an Empirical Demonstration''} for consideration as a
Research Paper in \emph{Artificial Intelligence in Medicine}.

\textbf{Why this paper.} The open actigraphy cohorts DEPRESJON, PSYKOSE and HYPERAKTIV
(Garcia-Ceja et al., 2018; Jakobsen et al., 2020; Hicks et al., 2021), now harmonised as
OBF-Psychiatric (Garcia-Ceja et al., \emph{Scientific Data}, 2025), are the most used public
benchmarks for machine-learning detection of mental disorders from wearable activity data.
Accuracies reported on them have climbed from the dataset authors' leave-one-patient-out 0.7 to
above 0.9. Our paper asks whether that reflects better models or weaker evaluation. Part~I is a
closed-scope systematic appraisal of the 28 published models on these cohorts, using the
PROBAST+AI risk-of-bias tool and the TRIPOD+AI reporting checklist. For each study we record
five things that can be checked against the data: whether training and test folds were split by
participant or by day, whether the model was tested on a second cohort, whether calibration was
reported, how many participants with the outcome were available per candidate predictor, and
whether code was released; we also record whether studies that combined DEPRESJON and PSYKOSE
noticed that the two share the same 32 control participants. Part~II re-evaluates the cohorts
with models fixed a priori, ten random seeds and participant-level bootstrap intervals, under
exactly the evaluation designs the appraisal scores as high and low risk.

\textbf{Main findings.} Splitting by day rather than by participant inflates AUROC by up to 0.17
and turns a chance-level ADHD classifier into an apparently useful one. Applying a frozen model to
a second cohort loses up to 0.21 AUROC in one direction, with collapse of calibration. Calibration
is poor even under honest validation, is not repairable from the training cohort alone, and is
never reported in the literature on these cohorts. Honest performance is stable across feature
subsets, resampling and cross-validation schemes. Every published claim above the honest band
coincides with a day-level or unstated splitting unit.

\textbf{Fit to the journal.} This is a methodological contribution to the evaluation of AI
prediction models in medicine. It provides reference baselines for widely used clinical
benchmarks and is itself reported according to the TRIPOD+AI guideline. All analysis code, run
manifests and generated tables are public, and the figures were regenerated independently on
Google Colab from the repository with identical results. The supplementary material contains the
TRIPOD+AI checklist, feature definitions, identification sources and the study list.

\textbf{Declarations.} The manuscript is original, is not under consideration elsewhere, and has
not been published in any form. Both authors approved the submission. No funding was received. We
declare no conflicts of interest. The study is a secondary analysis of publicly available,
de-identified datasets that their custodians released for research use, cited above and in the
manuscript. We wish to publish under the subscription (non-open-access) option.

\textbf{Suggested reviewers.} Because the contribution concerns the evaluation methodology of
machine-learning models, reviewers with expertise in machine learning for health data, model
validation and calibration, and reproducibility of AI systems are most appropriate; clinical
expertise in psychiatry is helpful but not required to assess the work. We have no exclusions.

Thank you for considering our manuscript.

\closing{Yours sincerely,}
\end{letter}
\end{document}
